\section{Análisis de resultados}

\subsection{Diseño y fabricación}
\subsubsection{Articulaciónes desmontables}

Con base en los lineamientos de tolerancias del apartado~\ref{Tolerancias}, se diseñó y fabricó una articulación desmontable para el brazo robótico. El objetivo fue verificar encastres, juego angular y comportamiento frente a rozamiento en las superficies de contacto.

\vspace{1em}

Inicialmente se construyó con descartes de MDF cortados a láser y herrajes M3. La pieza se sometió a esfuerzos de fricción, corte y rotación para observar desgaste, holguras y deformaciones. El prototipo permitió validar el principio mecánico, pero evidenció pérdida de tolerancia en contactos MDF-MDF bajo uso repetido. 

\begin{figure}[H]
  \centering
  \includegraphics[width=0.7\textwidth]{anexos/Articulacion/V2.png}
  \caption{Sistema de articulación desmontable. Fuente: elaboración própia}\label{fig:articulacionV2}
\end{figure}




Posteriormente se reemplazaron las superficies de contacto por insertos impresos en PLA, configurando una interfaz PLA–PLA en las zonas con rozamiento. Se ajustaron tolerancias según las galgas obtenidas, reduciendo el juego en el conjunto ensamblado y mejorando la repetibilidad. Se documentaron espesores, holguras objetivo y par de apriete de tornillería para su reproducción. Véase figura \ref{fig:corte_pernoV3}



\begin{figure}[H]
  \centering
  \includegraphics[width=0.7\textwidth]{anexos/Articulacion/Perno V3.png}
  \caption{Corte transversal perno --- versión mejorada. Fuente: elaboración própia}\label{fig:corte_pernoV3}
\end{figure}



Se identifican con verde las superficie de contacto que friccionan contra el perno azúl, las cualés tienen un ajuste deslizante con el perno. En amarillo se muestra la pieza que actúa como tope para el seguro del perno. En el modelo anterior se utilizaba una sola pieza que hacía esta función, pero al utilizar este tope (con un ajuste apretado) permite regular la fuerza de ajuste y juego de la articulación. El seguro que se encuentra entre el tope y el perno, este permite montar y desmontar la articulación de manera sencilla.


\begin{figure}[H]
  \centering
  \includegraphics[width=0.3\textwidth]{anexos/CAD/despieceArticulacion.png}
  \caption{Despiece de articulación desmontable. Fuente: elaboración própia}\label{fig:despiece_articulacion}
\end{figure}


\newpage

\subsubsection{Brazo robótico articulado}


\paragraph*{\textbf{Muñeca}} --- \\

La muñeca está compuesta por una estructura principal de 5 piezas de MDF (véase la figura~\ref{fig:brazoMuneca}). Alberga al servomotor de la misma, junto con el electroimán, y la articulación que la une al eslabón anterior.

\begin{figure}[H]
  \centering
  \includegraphics[width=0.3\textwidth]{anexos/CAD/brazoMuneca.png}
  \caption{Muñeca de brazo robótico. Fuente: elaboración própia}\label{fig:brazoMuneca}
\end{figure}


\paragraph*{\textbf{Eslabones}} --- \\

Los 2 eslabones se componen de 6 piezas de MDF pegadas entre sí para dar soporte estructural al servomotor y articulación desmotable respectivamente. Poseen encastres para el siguiente efector y articulación, además, como se ve en la figura~\ref{fig:brazoEslabon}, el eslabón inferior o principal, posee soportes exteriores para tornillos, los cuales trabajan en conjunto con el hombro para albergar gomas elásticas o un resorte, el cual se encargará de brindar asistencia mecánica en posiciones extendidas para la articulación inferior, la cual es la encargada de soportar la mayor carga. 


\begin{figure}[H]
  \centering
  \includegraphics[width=0.3\textwidth]{anexos/CAD/brazoEslabon.png}
  \caption{Eslabón de brazo robótico. Fuente: elaboración própia}\label{fig:brazoEslabon}
\end{figure}

\paragraph*{\textbf{Hombro}} --- \\

Este componente, como se puede observar en la figura~\ref{fig:brazoHombro}, es una pieza completmanete estructural. Sus funciones son brindar estabilidad estructural a los eslabones y a la muñeca, junto con el electroimán, permitir la rotación del brazo al acoplarse al eje del servomotor de la base, y aliviar cargas axiales en el mismo al apoyarse en el anillo de soporte. Además el hombro se encarga de contener a la suspensión de la articulación inferior, encargada de dar asistencia en posiciones de extensión máxima o cercanas, y distribuir el esfuerzo del servomotor más uniformemente en el rango de movimiento.

\begin{figure}[H]
  \centering
  \includegraphics[width=0.3\textwidth]{anexos/CAD/brazoHombro.png}
  \caption{Hombro de brazo robótico. Fuente: elaboración própia}\label{fig:brazoHombro}
\end{figure}


\paragraph*{\textbf{Base}} --- \\

La base del brazo se encarga de mantener estático al sistema durante operación (véase figura~\ref{fig:brazoBase}), a la vez de albergar el servomotor que se encarga de la rotación del brazo sobre su própio eje. Este se encuentra en la parte inferior de la base, y se encuentra asistido por un anillo que hace leve contacto entre la base y el hombro, lo que permite alivitar cargas axiales y momentos no deseados. 

\vspace{1em}

La base posee patas que se extienden lejos del centro de masa del sistema para brindar aproyo sobre cualquier extensión del brazo, sin que este llege a tambalearse. Además, la base cuenta con una mesa donde se alojan los componentes electrónicos, además del tablero. 

\begin{figure}[H]lidad 
  \centering
  \includegraphics[width=0.6\textwidth]{anexos/CAD/brazoBase.png}
  \caption{Base de brazo robótico. Fuente: elaboración própia}\label{fig:brazoBase}
\end{figure}


\paragraph*{\textbf{Tablero}} --- \\

Se incorpora un tablero desmontable el cual incorpora los componentes electrónicos para la comunicación del usuario, además de alojar conectores tanto para el controlador espejo como para la fuente de alimentación del sistema. 


\begin{figure}[H]
  \centering
  \includegraphics[width=0.3\textwidth]{anexos/CAD/brasoTablero.png}
  \caption{Tablero de control de brazo robótico. Fuente: elaboración própia}
\end{figure}







\subsubsection{Electroimán}
Para realizar la carcasa del electroimán se realiza uso de un torno manual. Hasta tener las medidas concretadas, se dejan un excedente en espesor y longitud para realizar pruebas durante el proceso de prototipado. Cuando se obtengan el resto de valores se dará lugar a una finalización de esta pieza por motivos de reducción de peso.

\vspace{1em}

La bobina del electroimán se construyó utilizando alambre de cobre esmaltado de 0,5 mm de diámetro, enrollado a lo largo de aproximadamente 5 metros para generar el campo magnético necesario. Como núcleo ferromagnético se empleó un tornillo de acero, seleccionado por su alta permeabilidad magnética y disponibilidad. 


\begin{figure}[H]
  \centering
  \includegraphics[width=0.5\textwidth]{anexos/Torno/Pieza final.jpg}
  \caption{Pieza fabricada. Fuente: elaboración própia}
\end{figure}

El electroimán cuenta con 5 partes principales, 4 de ellas se ven representadas en la figura~\ref{fig:despiece_iman}. El componente principal del electroimán es la la bobina de alambre de cobre esmaltado, esta es contenida por el carretel. Este se compone por doss piezas impresas en 3D que encastran entre sí. Esto se hizo así para evitar soportes durante el proceso de impresión. El carretel con la bobina es colocado dentro de la carcasa, el cual es asegurado por un anillo creado a partir de un segmento de barra de silicona, el cual evita que el carretel con la bobina salgan de la carcasa. Por último, se coloca un tornillo M10 o M9 en el centro que sirve tanto de núcleo para el electroimán, tanto como método de sujeción a la muñeca. 

\begin{figure}[H]
  \centering
  \includegraphics[width=0.3\textwidth]{anexos/CAD/despieceIman.png}
  \caption{Despiece de electroimán. Fuente: elaboración própia}\label{fig:despiece_iman}
\end{figure}




\subsection{Fuente de alimentación}

En lo referente a la electrónica de potencia, se utilizó una fuente ATX de una computadora en desuso para alimentar los motores del sistema. Esta fuente fue seleccionada debido a su capacidad para suministrar líneas estables de 12 V y 5 V, necesarias para los distintos actuadores.

\vspace{1em}

Además, se diseñó y cortó en MDF una carcasa para alojar y proteger la fuente, permitiendo un montaje seguro y una mejor organización del cableado dentro del proyecto, la misma se enciende por un interruptor ubicado en la parte delantera de la carcasa.


\begin{figure}[H]
  \centering
  \includegraphics[width=0.5\textwidth]{anexos/CAD/fuenteAlim.png}
  \caption{Modelo CAD de fuente de alimentación. Fuente: elaboración própia}
\end{figure}


\subsubsection{Controlador espejo}

Se realiza el diseño del controlador espejo tomando en consideración las tolerancias especuladas. Se elaboran bosquejos y se diseñan partes individuales para verificar calces y encastres, tomando como inspiración el anexo \ref{fig:insp.impreso}.

\vspace{1em}

Este controlador incorpora potenciometros para realizar la lectura de los ángulos de cada articulación, replicando la distribución de actuadores en el brazo real. A través de estos se puede realizar el control fluido del mismo, sin la necesidad de una programación extensiva de movimientos que pueden llegar a ser complicados de trazar sin cálculos complejos.

\vspace{1em}

Utilizando una filosofía de materiales híbrida, y luego de varias iteraciones, se llega al siguiente diseño: Un enfoque modular, con una estructura en MDF y con carcazas para protección y estética en PLA. La parte estructural está puramente dedicada con MDF por su velocidad de fabricación/prototipado, mientras que se utilzó PLA para carcazas por ser electricamente aislante, y su versatilidad al poder implementar impresión 3D como método de fabricación. La carcasa de la muñeca se encarga de albergar los botones para el electroimán y grabación.  

\begin{figure}[H]
  \centering
  \includegraphics[width=0.8\textwidth]{anexos/CAD/despieceEspejo.png}
  \caption{Despiece del controlador espejo. Fuente: elaboración própia}
\end{figure}

Los potenciometros se encuentran conectados en paralelo a 5V y tierra, mientras que el wiper se lo toma como referencia analógica para la lectura de los ángulos, además de conectar los pulsadores a tierra para la lectura de entradas de comando de electroimán y grabación utilizando las resistencias pull up internas del microcontrolador. Se crean conectores hembra para facilitar la manipulación de las conexiones eléctricas, tanto para la alimentación del controlador, como para salidas analógicas de los potenciomestros y las salidas digitales de los pulsadores.



\subsection{Electrónica de control}

% TODO: Redactar {Electrónica de control}

\begin{figure}[H]
  \centering
  \includegraphics[width=0.8\textwidth]{anexos/circuito/tinkerCAD.png}
  \caption{Diagrama de conexiones eléctricas. Fuente: elaboración própia}
\end{figure}




\subsection{Código microcontrolador}

% TODO: redactar Código microcontrolador

\begin{figure}[H]
  \centering
  \includegraphics[width=0.7\textwidth]{anexos/Codigo/Flujo.png}
  \caption{Diagrama de flujo del código en C. Fuente: elaboración própia}
\end{figure}






\subsection{Interfaz gráfica}
\subsubsection{Selección y justificación de la interfaz gráfica de control}

Se adoptó una aplicación de escritorio como interfaz humano-máquina (HMI) con el objetivo de facilitar el proceso de prototipado y permitir la incorporación rápida de nuevas funcionalidades. Este enfoque proporciona una mayor flexibilidad frente a los paneles físicos, evitando modificaciones en el hardware cada vez que se requieren ajustes en el sistema. Asimismo, la interfaz gráfica permite integrar visualización en tiempo real, ajustes finos de parámetros operativos y funciones avanzadas como la automatización de secuencias.

\subsubsection{Arquitectura distribuida del sistema}

El sistema se diseñó mediante una arquitectura distribuida en la cual la aplicación ejecuta las tareas de control de alto nivel, mientras que el microcontrolador gestiona el control de tiempo real de los actuadores. La comunicación entre ambos módulos se implementó mediante un enlace serial UART utilizando un protocolo textual simplificado. Esta separación permite mantener un comportamiento determinístico en el movimiento de los servomotores, al mismo tiempo que se preserva la flexibilidad y escalabilidad de la capa superior del sistema.

\subsubsection{Modos de operación implementados}

Se definieron dos modos principales de operación:
\begin{itemize}
    \item \textbf{Modo de control manual (modo espejo):} el robot replica directamente los movimientos ejecutados mediante el controlador físico.
    \item \textbf{Modo de control por aplicación:} las consignas de movimiento se generan desde la interfaz gráfica, lo que permite una mayor precisión y la posibilidad de automatizar secuencias.
\end{itemize}

La elección del modo se gestiona mediante lógica interna del sistema, garantizando transiciones suaves y sin interrupciones bruscas.

\subsubsection{Procedimiento de grabación y reproducción de trayectorias}

Se implementó un sistema de muestreo periódico de los ángulos articulares para registrar secuencias de movimiento. Las trayectorias se almacenan en estructuras ordenadas que preservan tanto el valor de cada articulación como el intervalo de muestreo. Para la reproducción, el sistema ejecuta las consignas respetando la cadencia temporal original, lo cual permite replicar movimientos de forma aproximada y consistente.

\subsubsection{Estrategias de seguridad aplicadas}

Con el fin de garantizar la integridad del sistema y del operador, se incorporaron diversas estrategias de seguridad en la capa de software. Entre ellas se incluyen:
\begin{itemize}
    \item Establecimiento de límites de operación por software para cada articulación.
    \item Validación de comandos provenientes de la aplicación para evitar consignas incoherentes.
    \item Implementación de transiciones progresivas entre modos de operación.
    \item Generación de rampas de movimiento al iniciar trayectorias automáticas, reduciendo esfuerzos mecánicos y evitando cambios bruscos.
\end{itemize}

\subsubsection{Integración entre la HMI y el microcontrolador}

La aplicación de escritorio envía comandos en un formato estructurado que el microcontrolador interpreta y ejecuta en tiempo real. De forma complementaria, el microcontrolador retorna información del estado del sistema, la cual es procesada y presentada visualmente al operador. Este mecanismo bidireccional garantiza coherencia entre la interfaz y el estado físico del robot, asegurando una operación confiable y sincronizada.


% TODO: Modelo CAD con todo incorporado

% TODO: Modelo fisico del brazo
% TODO: Modelo fisico del controlador espejo
% TODO: Modelo fisico de la fuente de alimentacions

% TODO: Diagrama de conexiones finales en TinkerCAD
% TODO: Circuito real

% TODO: Funcionamiento del código del arduino
% TODO: Funcionamiento de la interfaz de usuario

% TODO: Desempeño del dispositivo en modo remoto
% TODO: Desempeño del dispositivo en modo manual
% TODO: Desempeño del dispositivo en modo grabado

% TODO: Comparacion de los objetivos